\documentclass[11pt]{article}
\usepackage[margin=1.15in]{geometry}
\usepackage{amsmath, amssymb, amsthm}
\usepackage{booktabs}
\usepackage{xcolor}
\definecolor{linknavy}{RGB}{25, 55, 125}
\usepackage[colorlinks=true, linkcolor=linknavy, citecolor=linknavy, urlcolor=linknavy]{hyperref}
\newtheorem{corollary}{Corollary}
\newcommand{\C}{\mathbb{C}}

\title{The Collateral Impact of the Three-Dimensional Jacobian Counterexample:\\ a Source-Verified Cascade}
\author{Felipe Santiba\~nez-Leal\\
\small CAOS Research program, \texttt{github.com/fsantibanezleal/CAOS\_RESEARCH}}
\date{2026-07-21 (companion document, v0.01)}

\begin{document}
\maketitle

\begin{abstract}
On 2026-07-19, Levent Alp\"oge announced an explicit counterexample to the Jacobian conjecture
in dimension 3, found working with the language model Claude Fable (Anthropic). The credit for
every consequence documented here belongs to that single map. This companion records what else
that map decides: through implication chains published over five decades, and re-verified
against their primary sources by an experiment-history program (records EXP-001, EXP-016,
EXP-018), the counterexample resolves or partially resolves the Mathieu conjecture, the Dixmier
conjecture, the Poisson conjecture, the Gaussian moments conjecture, Zhao's vanishing
conjecture, the Image conjecture, and the symmetric (gradient) form of the Jacobian conjecture.
We state each corollary with its chain and its reference, list precisely what remains open, and
describe the verification method, whose value is exactly that a one-line certificate could be
turned into a reliable status map of a whole neighborhood of problems within two days.
\end{abstract}

\section{The event, and the discipline applied to it}

The counterexample is the polynomial map $F : \C^3 \to \C^3$, $u = 1 + xy$,
\[
F = \bigl(u^3 z + y^2 u (4 + 3xy),\;\; y + 3x u^2 z + 3x y^2 (4 + 3xy),\;\; 2x - 3x^2 y - x^3 z\bigr),
\]
with $\det JF = -2$ everywhere and
$F(0, 0, -\tfrac14) = F(1, -\tfrac32, \tfrac{13}{2}) = F(-1, \tfrac32, \tfrac{13}{2})$
\cite{Alpoge2026}. Our program validated it independently in exact rational arithmetic and
computed the full fiber over the repeated value: exactly the three announced points (record
EXP-001). Every consequence below is a corollary of this validation plus one or more published
implications; for each implication the primary record was located and its direction confirmed
before the corollary was asserted (records EXP-016 and EXP-018). Nothing here is original
mathematics of ours beyond the verification work and the bookkeeping discipline.

\section{The cascade}

\begin{corollary}[Smale's 16th problem]
The Jacobian conjecture is false for every $N \ge 3$; the two-variable case remains open.
\end{corollary}

\begin{corollary}[Mathieu conjecture, 1997]
Mathieu proved that his conjecture for $SU(N)$ implies $JC(N)$ \cite{Mathieu1997, Zwart2025};
Duistermaat and van der Kallen proved the abelian case \cite{DvdK1998}. By contraposition the
Mathieu conjecture is \textbf{false for $SU(N)$, $N \ge 3$}. It remains open for $SU(2)$ and
for the other compact connected groups.
\end{corollary}

\begin{corollary}[Dixmier conjecture, 1968]
$JC(2n)$ implies $\mathrm{Dixmier}(n)$, and $\mathrm{Dixmier}(n)$ implies $JC(n)$
\cite{BKK2007, Tsuchimoto2005}. Since $JC(4)$ is false, $\mathrm{Dixmier}(2)$ is false: the
\textbf{full Dixmier conjecture is false}. The original rank-1 case remains open.
\end{corollary}

\begin{corollary}[Poisson conjecture, 2007]
The Jacobian, Dixmier and Poisson conjectures are equivalent in any characteristic
\cite{AvdE2007}. Hence the \textbf{full Poisson conjecture is false}; the minimal failing
dimension is open.
\end{corollary}

\begin{corollary}[Gaussian moments conjecture]
GMC implies JC \cite{DEZ2017}; hence \textbf{GMC is false} at some finite dimension.
\end{corollary}

\begin{corollary}[Zhao's vanishing conjecture]
The vanishing conjecture for Hessian-nilpotent polynomials is equivalent to JC
\cite{Zhao2004, Zhao2007}; hence it is \textbf{false}: there exists a Hessian-nilpotent
quartic $P$ with $\Delta^m P^m = 0$ for all $m$ but $\Delta^m P^{m+1} \ne 0$ for infinitely
many $m$. An explicit witness is a target of the program.
\end{corollary}

\begin{corollary}[Image conjecture]
The Image conjecture implies the vanishing conjecture \cite{vdE2010}; hence it is
\textbf{false in some dimension}.
\end{corollary}

\begin{corollary}[Symmetric / gradient Jacobian conjecture]
JC reduces to cubic homogeneous maps with symmetric Jacobian \cite{dBvdE2005}; with the Zhao
equivalence this yields a dimension in which a \textbf{gradient Keller map is non-invertible}.
\end{corollary}

\section{What remains open}

The two-variable Jacobian conjecture; Dixmier at rank 1; Mathieu for $SU(2)$ and the remaining
compact groups; the minimal failing dimensions of the Poisson, moments, vanishing and Image
statements; and every explicit witness beyond dimension bookkeeping (the program's queued
target: extracting a failing Hessian-nilpotent quartic from the counterexample family).
Excluded deliberately from this list: the Markus--Yamabe conjecture (refuted independently in
1997, not a corollary of this event) and statements whose implication direction does not
propagate the falsity of JC.

\section{Method note}

The verification program is an open experiment-history repository: every experiment declares a
falsifiable hypothesis before running, certifies with exact arithmetic (never floats), closes
with a persisted verdict, and keeps refuted attempts in the record. For this document the
relevant records are EXP-001 (exact validation, full fiber), EXP-016 and EXP-018 (each chain
located in its primary source and its direction confirmed). The general lesson we would offer:
when a machine-checkable certificate lands, the fastest reliable way to know what it changes is
not commentary but a systematic, source-verified sweep of the implication graph, and that sweep
is itself an experiment with a hypothesis, a method and a verdict.

\begin{thebibliography}{9}
\bibitem{Alpoge2026} L. Alp\"oge, announcement of the counterexample, X (@\_\_alpoge\_\_),
2026-07-20; found working with Claude Fable (Anthropic); question posed by Akhil.
\bibitem{Mathieu1997} O. Mathieu, Some conjectures about invariant theory and their
applications, 1997.
\bibitem{DvdK1998} J. J. Duistermaat, W. van der Kallen, Constant terms in powers of a Laurent
polynomial, \emph{Indag. Math.} 9 (1998), 221--231.
\bibitem{Zwart2025} K. Zwart, Mathieu's approach to the Jacobian Conjecture (review),
arXiv:2511.16561.
\bibitem{BKK2007} A. Belov-Kanel, M. Kontsevich, The Jacobian conjecture is stably equivalent
to the Dixmier conjecture, arXiv:math/0512171; \emph{Mosc. Math. J.} 7 (2007).
\bibitem{Tsuchimoto2005} Y. Tsuchimoto, Endomorphisms of Weyl algebra and $p$-curvatures,
\emph{Osaka J. Math.} 42 (2005), 435--452.
\bibitem{AvdE2007} P. K. Adjamagbo, A. van den Essen, On the equivalence of the Jacobian,
Dixmier and Poisson conjectures in any characteristic, arXiv:math/0608009; \emph{Acta Math.
Vietnam.} 32 (2007), 209--218.
\bibitem{DEZ2017} H. Derksen, A. van den Essen, W. Zhao, The Gaussian moments conjecture and
the Jacobian conjecture, \emph{Israel J. Math.} 219 (2017), 917--928; arXiv:1506.05192.
\bibitem{Zhao2004} W. Zhao, Hessian nilpotent polynomials and the Jacobian conjecture,
arXiv:math/0409534; \emph{Trans. Amer. Math. Soc.} 359 (2007).
\bibitem{Zhao2007} W. Zhao, A vanishing conjecture on differential operators with constant
coefficients, arXiv:0704.1691.
\bibitem{vdE2010} A. van den Essen, The amazing Image conjecture, arXiv:1006.5801.
\bibitem{dBvdE2005} M. de Bondt, A. van den Essen, A reduction of the Jacobian conjecture to
the symmetric case, \emph{Proc. Amer. Math. Soc.} 133 (2005), 2201--2205.
\end{thebibliography}

\end{document}
